\section{Related Work}
\label{sec:related}

\paragraph{Aggregation and task selection in multi-task evaluation.}
Multi-task results depend on which tasks are included and how they are combined.
\citet{dehghani2021lottery} show that rankings can change with task selection and discuss macro-averaging, arithmetic versus geometric means and rank-based aggregation; \citet{liang2023helm} leave a universal single-number aggregate explicitly unresolved and stress that it encodes value judgements.
These works vary the benchmark while the compared models stay as they are. We instead hold the benchmark and all rivals fixed and edit one model's scores.
They do not establish an ordering of rules by the cost of changing a ranking, and we do not read one into them.

\paragraph{Social choice and evaluation.}
\citet{lanctot2023vase} cast evaluation as voting, with models as alternatives and tasks as voters, so that comparisons need only within-task order. The Handbook of Computational Social Choice \citep{brandt2016handbook} supplies the voting-rule and Condorcet vocabulary we use \citep{zwicker2016voting}.
Our target-promotion problem is related to priced bribery, where a fixed set of voters is paid to change their votes \citep{faliszewski2016control}; it is a structured special case in which the target may only move up on the tasks it is allowed to edit, so complexity results for unrestricted bribery do not transfer automatically.
We make no claim of new algorithms for bribery or control.

\paragraph{Leaderboard incentives and strategic post-training.}
\citet{singh2025illusion} document private testing, selective disclosure and data-access asymmetries on Chatbot Arena, and find in a controlled training-mixture experiment that gains on the targeted evaluation need not transfer broadly.
\citet{chen2026incentives} analyse simultaneous developers' effort and a designer's tuning policy as an equilibrium problem, under assumptions on the score response to effort.
Our setting is different and narrower: a single target, fixed rivals, observed score matrices, and no equilibrium, dynamics or latent capability.
It measures the \emph{score improvement} that a ranking goal requires under a stated cap, not the training effort that would deliver it.

\paragraph{Benchmark-related training and contamination.}
\citet{zhou2023cheater} show that continued training on benchmark-related data raises benchmark scores, with possible losses elsewhere, and \citet{yang2023rephrased} show that rephrased test samples can evade common decontamination while still inflating scores.
These works motivate the intervention class we study; they do not provide a task-response or cost law, and they show that independent monotone gains with no side effects are an abstraction (Section~\ref{sec:limits}).
Score tables alone neither identify nor exclude contamination, and nothing here accuses any listed model.

\paragraph{Benchmarks and data used.}
We use the HELM MMLU release \citep{liang2023helm,hendrycks2021mmlu} and the Open LLM Leaderboard, whose tasks include BBH \citep{suzgun2022bbh} and which is built on the Evaluation Harness \citep{sutawika2023harness}.
